\documentclass{scrartcl}
\usepackage[latin1]{inputenc}
\usepackage[T1]{fontenc}
\usepackage[ngerman]{babel}
\usepackage{amsmath}
\usepackage{amssymb}
\usepackage{icomma}
\usepackage{nicefrac}
%\usepackage[dvips]{graphicx}
%\usepackage{floatflt}
%\usepackage{enumitem}
%\usepackage{babel}
\usepackage{blindtext}
%\usepackage{showframe}
\usepackage{calc}
\usepackage{wrapfig}
\def\BILD{\rule{0.4\textwidth}{4cm}}

\usepackage{graphicx}
\usepackage{placeins}
\usepackage{multirow}
\usepackage{subfig}
\usepackage{url}

\renewcommand{\topfraction}{.85}
\renewcommand{\bottomfraction}{.7}
\renewcommand{\textfraction}{.15}
\renewcommand{\floatpagefraction}{.66}
\renewcommand{\dbltopfraction}{.66}
\renewcommand{\dblfloatpagefraction}{.66}
\setcounter{topnumber}{9}
\setcounter{bottomnumber}{9}
\setcounter{totalnumber}{20}
\setcounter{dbltopnumber}{9}
\setlength{\intextsep}{0cm plus1cm minus1cm}
\setlength{\parindent}{0cm}
\pdfminorversion = 5
\usepackage{setspace}
\onehalfspacing

\begin{document}
\title{Versuch 11: Atomspektren}

\date{\today}

\author{Gruppe 5a: Gia-Danh Lam, Nils Haldenwang}

\maketitle
\tableofcontents

\newpage
\section{Einleitung und theoretischer Hintergrund}
\label{sec:einleitung}

In einer Atomh�lle gebundene Elektronen k�nnen sich auf
unterschiedlichen Energieniveaus aufhalten. Durch Absorption eines
Photons der Energie $ E = h \cdot \nu $ kann das Elektron auf ein
h�heres Energieniveau angehoben werden. Dabei entspricht die Energie
des Photons genau der Energiedifferenz der beiden Niveaus. Umgekehrt
wird ein Photon entsprechender Energie emittiert, wenn ein Elektron
von einem h�heren Energieniveau auf ein niedrigeres f�llt. Da die
Energieniveaus diskret sind, m�ssen auch ihre Differenzen diskret
sein. Dies ist die Ursache daf�r, dass Emissionspektren von Atomen
stets aus einzelnen Linien bestehen. Man bezeichnet sie daher auch als
Linienspektren.

\subsection{Spektrum des Wasserstoffatoms}
\label{sec:spektr-des-wass}

F�r Wasserstoff erh�lt man als Gleichung f�r die Energie der einzelnen Schalen:

\begin{equation}
  W_{n}= R \cdot\frac{1}{n^2_1}
\label{eq:1}
\end{equation}

mit

\[R = -
\frac{e^4m}{2\hbar�}\frac{1}{(4\pi\epsilon_0)^2} =-
\frac{e^4m}{8h^2\epsilon_0^2}  \]

dabei ist $e = Elementarladung$, $m = Elektronenmasse$, $\hbar =
\frac{h}{2\pi}$ und $n = Hauptquantenzahl$. Die sog. Rydbergkonstante
hat den Wert $R = -13.605 eV$. Sie ist eine der am besten vermessenen
Naturkonstanten.

Damit erh�lt man f�r die Energiedifferenz Eines �bergangs von $n_2$ zu
$n_1$:

\begin{equation}
W_{n_2\mapsto n_1}=-R\left(\frac{1}{n_1^2}-\frac{1}{n_2^2}\right)
\label{eq:2}
\end{equation}

Mit Hilfe von $W=h\cdot \nu$, $\lambda= \frac{c}{\nu}$ und $R=
R_H\cdot h\cdot c$ ergibt sich f�r die Wellenl�nge des absorbierten/emittierten Photons:

\begin{equation}
\lambda = \frac{1}{R_H \left(\frac{1}{n_1^2}- \frac{1}{n_2^2}\right)} \text { mit } R_H=1,097\cdot 10^5 \frac{1}{cm}
\label{eq:3}
\end{equation}

Als \emph{Serien} bezeichnet man alle �berg�nge von beliebigen
Energieniveaus auf ein bestimmtes Zielniveau, also $n_1 = const.$ und
$n_2 \geq n_1 + 1$.

\begin{figure}%[htb!]
  \centering
    \begin{minipage}[c]{7.2 cm}
    \includegraphics[width=8cm]{pics/dummypic}
    \caption{Spektralserien des Wasserstoffs ( nicht alle Linien
      dargestellt ).}
    \label{fig:serien}
  \end{minipage}
  \begin{minipage}[c]{7.2 cm}
    \includegraphics[width=8cm]{pics/dummypic}
    \caption{Grotrian-Diagramm f�r Kalium}
    \label{fig:grotrian}
  \end{minipage}
\end{figure}

\section{Doubletts durch Spin-Bahn-Kopplung}
\label{sec:doubletts-durch-spin}

Alkalimetalle wie Kalium besitzen �hnlich dem Wasserstoff nur ein
s-Valenzelektron im elektronischen Grundzustand. Dadurch ergibt sich
eine gewisse �hnlichkeit der Spektren. Markante Unterschiede entstehen
durch die sog. \emph{Spin-Bahn-Kopplung}. Spindrehimpuls und
Bahndrehimpuls des Elektrons treten dabei miteinander in
Wechselwirkung. Die beiden Momente k�nnen entweder parallel oder
entgegengesetzt orientiert sein, sodass sich f�r ein Niveau mit
Bahndrehimpuls $l$ eine Aufspaltung des Gesamtdrehimpulses $j$ in $j =
l + \frac{1}{2}$ und $j = l - \frac{1}{2}$ ergibt. Es existieren also
im Linienspektrum sehr nah beieinander liegende Linien, die mit
ausreichendem Aufl�sungsverm�gen auch sichtbar gemacht werden k�nnen.
Im in Abbildung \ref{fig:grotrian} dargestellten Ausschnitt aus dem
Grotrian-Diagramm f�r Kalium sind zwei Doubletts erkennbar f�r die
�berg�nge $5^2P_{1 \pm \nicefrac{1}{2}} \mapsto 4^2S_{\nicefrac{1}{2}}$ und
$4^2P_{1 \pm \nicefrac{1}{2}} \mapsto 4^2S_{\nicefrac{1}{2}}$.

Die Notation f�r die Termschemata ist dabei die
folgende, wobei die Bedeutung der Symbole in Tabelle \ref{tab:
  Quantenzahlen} dargestellt sind:

\begin{equation}
n ^{(2s+1)} L_{j}
\label{eq:Notation}
\end{equation}

\begin{table}%[htb!]
  \centering
  \begin{tabular}{| c | l |}
    \hline
    Symbol & Bedeutung \\
    \hline
    n	&	Hauptquantenzahl ( Schalen ) \\
    L	&	Drehimpulsquantenzahl $\widehat=\: l $ ( S,P,D,F,... ) \\
    l	&	Bahndrehimpulsquantenzahl	\\
    s	&	Spinquantenzahl $ s = \pm \frac{1}{2}$\\
    j	& Gesamtdrehimpulsquantenzahl $j = l + s$ \\
    \hline
  \end{tabular}
  \caption{Quantenzahlen f�r die Notation der Termsymbole.}
  \label{tab: Quantenzahlen}
\end{table}

Das Natriumspektrum wiederum unterscheidet sich deutlich von dem des
Wasserstoffs. Im sichtbaren Bereich enth�lt es bis auf einige sehr
intensit�tsschwache Ausnahmen nur die gelbe Na-Linie. Bei genauerer
Betrachtung zeigt sich, dass diese tats�chlich aber ebenfalls ein
Dublette ist mit den in Tabelle \ref{tab:na-doublette} dargestellten
Wellenl�ngen und �berg�ngen.

\begin{table}%[htb!]
  \centering
  \begin{tabular}{| c | c |}
    \hline
    �bergang & Wellenl�nge \\
    \hline
    $3 ^2P_{3/2} \mapsto 3 ^2S_{1/2}$ & 588.995nm\\
    \hline
    $3 ^2P_{1/2} \mapsto 3 ^2S_{1/2}$ & 589.592nm\\
    \hline
  \end{tabular}
  \caption{Natrium Doublette}
  \label{tab:na-doublette}
\end{table}

\end{document}



